\subsection{Restricted formal power series}

DEFINITION 1. A commutative topological ring A is said to be linearly topologized (and its topology is said to be linear) if there exists a fundamental system B of neighbourhoods of 0 consisting of ideals of A.

Note that in such a ring, the ideals \(\mathfrak{S} \in \mathcal{B}\) are open and closed (General Topology, Chapter III, § 2, no. 1, Corollary to Proposition 4). For all \(\mathfrak{S} \in \mathcal{B}\) , the quotient topological ring A/ \(\mathfrak{S}\) is then discrete; for \(\mathfrak{S} \in \mathcal{B}\) , \(\mathfrak{S}' \in \mathcal{B}\) , \(\mathfrak{S}' \subset \mathfrak{S}\) , let

\[
h _ {\mathfrak {S S} ^ {\prime}} \colon \mathrm{A} / \mathfrak {S} ^ {\prime} \rightarrow \mathrm{A} / \mathfrak {S}
\]

be the canonical mapping. We know (General Topology, Chapter III, § 7, no. 3) that \((\mathbf{A} / \mathfrak{S}, h_{\mathfrak{S}\mathfrak{S}'})\) is an inverse system of discrete rings (relative to the indexing set \(\mathcal{B}\) which is ordered by \(\supset\) and directed), whose inverse limit is a complete Hausdorff linearly topologized ring \(\tilde{\mathbf{A}}\) ; further (loc. cit., Proposition 2), a strict morphism \(i: \mathbf{A} \to \tilde{\mathbf{A}}\) is defined, whose kernel is the closure of \(\{0\}\) in \(\mathbf{A}\) and whose image is everywhere dense in \(\tilde{\mathbf{A}}\) , so that \(\tilde{\mathbf{A}}\) is canonically identified with the Hausdorff completion of \(\mathbf{A}\) .

DEFINITION 2. Given a commutative topological ring A, a formal power series

\[
\mathbf {T} = \sum_ {(n _ {i})} c _ {n _ {1} n _ {2} \dots n _ {p}} \mathbf {X} _ {1} ^ {n _ {1}} \mathbf {X} _ {2} ^ {n _ {2}} \dots \mathbf {X} _ {p} ^ {n _ {p}}
\]

in the ring \(\mathbf{A}[[\mathbf{X}_1,\ldots ,\mathbf{X}_p]]\) is called restricted if, for every neighbourhood \(\mathbf{V}\) of 0 in A, there is only a finite number of coefficients \(c_{n_1n_2\dots n_p}\) not belonging to \(\mathbf{V}\) (in other words, the family \((c_{n_1n_2\dots n_p})\) tends to 0 in A with respect to the filter of complements of finite subsets of \(\mathbf{N}^p\) ).

If A is linearly topologized, the restricted formal power series in \(A[[X_{1},\ldots,X_{p}]]\) form a subring of \(A[[X_{1},\ldots,X_{p}]]\) , denoted by \(A\{X_{1},\ldots,X_{p}\}\) : for if \(T=\sum_{(n_{i})}c_{n_{1}\ldots n_{p}}X_{1}^{n_{1}}\ldots X_{p}^{n_{p}}\) , \(T'=\sum_{(n_{i})}c_{n_{1}\ldots n_{p}}^{'}X_{1}^{n_{1}}\ldots X_{p}^{n_{p}}\) are two restricted formal power series and \(\mathfrak{S}\) a neighbourhood of 0 in A which is an ideal of A, there exists an integer \(m\) such that \(c_{n_1\ldots n_p} \in \mathfrak{S}\) and \(c_{n_1\ldots n_p}' \in \mathfrak{S}\) for every system \((n_1, \ldots, n_p)\) such that \(n_k \geqslant m\) for at least one index \(k\) ; now, if

\[
\mathrm{T} ^ {\prime \prime} = \mathrm{TT} ^ {\prime} = \sum_ {(n _ {i})} c _ {n _ {1} \dots n _ {p}} ^ {\prime \prime} \mathrm{X} _ {1} ^ {n _ {1}} \dots \mathrm{X} _ {p} ^ {n _ {p}},
\]

then \(c_{n_1 \ldots n_p}'' = \sum c_{r_1 \ldots r_p} c_{s_1 \ldots s_p}'\) for all systems \((r_k), (s_k)\) such that \(r_k + s_k = n_k\) for \(1 \leqslant k \leqslant p\) ; we conclude that if \(n_k \geqslant 2m\) , then \(r_k \geqslant m\) or \(s_k \geqslant m\) and hence, since \(\mathfrak{S}\) is an ideal, \(c_{n_1 \ldots n_p}'' \in \mathfrak{S}\) so long as \(n_k \geqslant 2m\) for at least one \(k\) , which establishes our assertion. Moreover, every derivative \(\partial T / \partial X_i\) ( \(1 \leqslant i \leqslant p\) ) of a restricted formal power series is restricted, as follows immediately from the definition and the fact that the neighbourhoods \(\mathfrak{S} \in \mathcal{B}\) are additive subgroups of A.

If A is discrete, the ring of restricted formal power series is just the polynomial ring \(A[X_{1},\ldots,X_{p}]\) .

Let us always assume that A is linearly topologized and let \(\mathcal{B}\) be a fundamental system of neighbourhoods of 0 in A consisting of ideals of A; for all \(\mathfrak{S} \in \mathcal{B}\) , let \(p_{\mathfrak{S}}: A \to A/\mathfrak{S}\) be the canonical homomorphism. By definition, for every restricted formal power series \(T \in A\{X_{1}, \ldots, X_{p}\}\) ,

\[
\bar {p} _ {\Im} (\mathrm{T}) \in (\mathrm{A} / \Im) [ \mathrm{X} _ {1}, \dots , \mathrm{X} _ {p} ].
\]

Clearly

\[
((\mathbf {A} / \mathfrak {S}) [ \mathbf {X _ {1}}, \dots , \mathbf {X _ {p}} ], \bar {h} _ {\mathfrak {S} \mathfrak {S} ^ {\prime}})
\]

is an inverse system of rings (relative to the directed indexing set \(\mathcal{B}\) ) and \((\bar{p}_{3})\) is an inverse system of homomorphisms \(\mathrm{A}\{\mathrm{X}_1,\dots ,\mathrm{X}_p\} \to (\mathrm{A} / \Im)[\mathrm{X}_1,\dots ,\mathrm{X}_p]\) ; as every polynomial is a restricted formal power series, \(\bar{p}_{3}\) is surjective; its kernel \(\mathrm{N}_{3}\) is the ideal of \(\mathrm{A}\{\mathrm{X}_1,\dots ,\mathrm{X}_p\}\) consisting of the restricted formal power series all of whose coefficients belong to \(\Im\) ; we shall give \(\mathrm{A}\{\mathrm{X}_1,\dots ,\mathrm{X}_p\}\) the (linear) topology for which the \(\mathrm{N}_{3}\) (for \(\Im \in \mathcal{B}\) ) form a fundamental system of neighbourhoods of 0 (a topology which obviously depends only on that on A). Then it follows from General Topology, Chapter III, § 7, no. 3, Proposition 2 that

\[
\pi = \varprojlim_ {\mathfrak {S}} \bar {p} _ {\mathfrak {S}}: \mathrm{A} \{\mathrm{X} _ {1}, \dots , \mathrm{X} _ {p} \} \rightarrow \varprojlim_ {\mathfrak {S}} (\mathrm{A} / \mathfrak {S}) [ \mathrm{X} _ {1}, \dots , \mathrm{X} _ {p} ]\tag{3}
\]

is a strict morphism whose kernel is the closure of \(\{0\}\) in \(\mathbf{A}\{\mathbf{X}_1,\ldots ,\mathbf{X}_p\}\) and whose image is dense in

\[
\mathrm{A} ^ {\prime} = \lim _ {\leftarrow \mathfrak {S}} (\mathrm{A} / \mathfrak {S}) [ \mathrm{X} _ {1}, \dots , \mathrm{X} _ {p} ].
\]

PROPOSITION 3. If the linearly topologized commutative ring A is Hausdorff and complete, the canonical homomorphism \(\pi\) is a topological ring isomorphism.

For all \((n_1, \ldots, n_p) \in \mathbf{N}^p\) and all \(\mathfrak{S} \in \mathcal{B}\) , let \(\phi_{n_1 \ldots n_p}^{\mathfrak{S}}\) be the mapping \((\mathrm{A} / \mathfrak{S})[\mathrm{X}_1, \ldots, \mathrm{X}_p] \to \mathrm{A} / \mathfrak{S}\) which maps every polynomial to the coefficient of \(\mathrm{X}_1^{n_1} \ldots \mathrm{X}_p^{n_p}\) in this polynomial; clearly the \(\phi_{n_1 \ldots n_p}^{\mathfrak{S}}\) form an inverse system of \((\mathrm{A} / \mathfrak{S})\) -module homomorphisms (relative to the ordered set \(\mathcal{B}\) ) and, as \(\mathbf{A}\) is canonically identified with \(\varprojlim_{\mathfrak{S}} (\mathrm{A} / \mathfrak{S})\) by hypothesis, \(\phi_{n_1 \ldots n_p} = \varprojlim_{\mathfrak{S}} \phi_{n_1 \ldots n_p}^{\mathfrak{S}}\) is a continuous A-homomorphism from \(\mathrm{A}'\) to \(\mathbf{A}\) . For every element \(S = (S_{\mathfrak{S}})_{\mathfrak{S} \in \mathcal{B}}\) of \(\mathrm{A}'\) , we shall see that the formal power series \(T = \sum_{(n_i)} \phi_{n_1 \ldots n_p}(S) X_1^{n_1} \ldots X_p^{n_p}\) is restricted and satisfies \(\pi(T) = S\) . For all \(\mathfrak{S} \in \mathcal{B}\) and all \(\mathfrak{S}' \in \mathcal{B}\) such that \(\mathfrak{S}' \subset \mathfrak{S}\) , the relation \(\phi_{n_1 \ldots n_p}^{\mathfrak{S}}(S_{\mathfrak{S}}) = 0\) implies

\[
\phi_ {n _ {1} \dots n _ {p}} ^ {\mathfrak {S} ^ {\prime}} (\mathrm{S} _ {\mathfrak {S} ^ {\prime}}) \in \mathfrak {S} / \mathfrak {S} ^ {\prime};
\]

as \(S_{\mathfrak{S}}\) is a polynomial, we see that \(\phi_{n_{1}\ldots n_{p}}(S) \in \mathfrak{S}\) except for those \((n_{1}, \ldots, n_{p})\) , finite in number, such that \(\phi_{n_{1}\ldots n_{p}}^{\mathfrak{S}}(S_{\mathfrak{S}}) \neq 0\) , which proves our first assertion; the second follows from the definitions. As A is Hausdorff, the intersection of the \(N_{\mathfrak{S}}\) reduces to 0 and hence \(\pi\) is bijective, which completes the proof, since \(\pi\) is a strict morphism.

PROPOSITION 4. Let A, B be two linearly topologized commutative rings, B being Hausdorff and complete, and \(u: \mathbf{A} \to \mathbf{B}\) a continuous homomorphism. For every family \(\mathbf{b} = (b_i)_{1 \leq i \leq p}\) of elements of B, there exists a unique continuous homomorphism

\[
\tilde {u} \colon \mathrm{A} \{\mathrm{X} _ {1}, \dots , \mathrm{X} _ {p} \} \rightarrow \mathrm{B}
\]

such that \(\tilde{u}(a) = u(a)\) for all \(a \in \mathbf{A}\) and \(\tilde{u}(\mathbf{X}_i) = b_i\) for \(1 \leqslant i \leqslant p\) .

There exists a unique homomorphism \(v \colon A[X_1, \ldots, X_p] \to B\) such that \(v(a) = u(a)\) for \(a \in A\) and \(v(X_i) = b_i\) for \(1 \leqslant i \leqslant p\) . Moreover, if \(\mathfrak{H}\) is a neighbourhood of 0 in \(B\) which is an ideal, \(\bar{u}^{-1}(\mathfrak{H}) = \mathfrak{I}\) is an ideal of \(A\) which is a neighbourhood of 0 and, for every polynomial \(P \in N_{\mathfrak{S}}\) , clearly \(v(P) \in \mathfrak{H}\) and hence \(v\) is continuous. As \(A[X_1, \ldots, X_p]\) is dense in \(A\{X_1, \ldots, X_p\}\) , the existence and uniqueness of \(\tilde{u}\) follow from General Topology, Chapter III, § 3, no. 3, Proposition 5 and the principle of extension of identities.

In the special case where \(\mathbf{A} = \mathbf{B}\) and \(u\) is the identity mapping we shall write \(f(b_{1},\ldots ,b_{p})\) or \(f(\mathbf{b})\) for the value of \(\tilde{u} (f)\) for every restricted formal power series \(f\in \mathrm{A}\{\mathbf{X}_1,\dots ,\mathbf{X}_p\}\) .

\textbf{Remarks}\par

\begin{enumerate}
\def\labelenumi{(\arabic{enumi})}
\item
  Proposition 4 proves that for every closed ideal \(\mathfrak{a}\) in a ring A which is assumed to be Hausdorff and complete, the relations \(b_{i} \in \mathfrak{a}\) for \(1 \leqslant i \leqslant p\) imply \(f(b_{1}, \ldots, b_{p}) \in \mathfrak{a}\) for every restricted formal power series \(f \in \mathrm{A}\{\mathrm{X}_{1}, \ldots, \mathrm{X}_{p}\}\) .
\item
  Suppose that \(\mathbf{A}\) is linearly topologized; let \(r\) be an integer such that
\end{enumerate}

\(1 \leqslant r \leqslant p\) and let the ring \(\mathrm{A}\{\mathrm{X}_1, \ldots, \mathrm{X}_r\}\) be given the topology defined above. Then the topological ring \(\mathrm{A}\{\mathrm{X}_1, \ldots, \mathrm{X}_p\}\) is identified with the ring of restricted formal power series

\[
(\mathrm{A} \{\mathrm{X} _ {1}, \dots , \mathrm{X} _ {r} \}) \{\mathrm{X} _ {r + 1}, \dots , \mathrm{X} _ {p} \}
\]

as follows immediately from the definitions.

\begin{enumerate}
\def\labelenumi{(\arabic{enumi})}
\setcounter{enumi}{2}
\tightlist
\item
  With the notation of Remark 2, suppose further that A is Hausdorff and complete and let us write every restricted formal power series \(f \in \mathrm{A}\{\mathbf{X}_1, \ldots, \mathbf{X}_p\}\) in the form
\end{enumerate}

\[
f = \sum_ {(n _ {i})} c _ {n _ {r + 1} \dots n _ {p}} (\mathrm{X} _ {1}, \dots , \mathrm{X} _ {r}) \mathrm{X} _ {r + 1} ^ {n _ {r + 1}} \dots \mathrm{X} _ {p} ^ {n _ {p}}
\]

where the \(c_{n_{r+1}\ldots n_p}\) are restricted formal power series. For every system \(\mathbf{x} = (x_1, \ldots, x_r)\) of elements of A, let

\[
b _ {n _ {r + 1} \dots n _ {p}} = c _ {n _ {r + 1} \dots n _ {p}} (x _ {1}, \dots , x _ {p}).
\]

It follows immediately from Remark 1 that \(\sum_{(n_i)} b_{n_r+1 \ldots n_p} X_{r+1}^{n_r+1} \ldots X_p^{n_p}\) is a restricted formal power series denoted by \(f(x_1, \ldots, x_r, X_{r+1}, \ldots, X_p)\) ; it is said to be obtained by substituting the \(x_i\) for the \(X_i\) for \(1 \leqslant i \leqslant r\) in \(f\) .

